\section{Buying the full-mesh finite-difference bounds}
\label{sec:ex-fullmesh-fixedbounds}

We prove the finite-volume estimate needed before any continuum extraction. The
important order is to pay for the bounded projected coefficient first, and only
then use the exact Wilson difference identity. The result is uniform over every
period at all sufficiently fine meshes. The proof has two parts: a finite
dimensional absorption estimate for one coefficient row, and a packing argument
which plants finitely many such rows into one large periodic box.

For a period $p$, mesh $N$, field list $v=(v_1,\ldots,v_n)$, and finite
list $l$ of slot-and-coordinate directions, write $K_{N,p;v,l}(x)$ for the
finite Wilson expectation of the product of the corresponding iterated
periodic difference insertions at the positions $x=(x_1,\ldots,x_n)$. The
normalization is $\Delta_{N,\alpha}F(x)=a_N^{-1}
\bigl(F(x+a_Ne_\alpha)-F(x)\bigr)$. Thus the theorem below concerns the unscaled
derivative kernel; the exact one-step identity used in its proof is
$F(x+a_Ne_\alpha)-F(x)=a_N\Delta_{N,\alpha}F(x)$.

\begin{lemma}[Absorption of the projected coefficient]\label{lem:ex-fullmesh-cap}
Let $I$ be a directed set of mesh and period indices, and let $X_i$ be a
finite-dimensional normed coefficient space. Suppose that for each $i$ there
are a projection $P_i:X_i\to X_i$, a linear map $L_i:X_i\to Y_i$, and a constant $Z>0$
such that
\[
   \|P_i z\|\le Z\|L_i(P_i z)\|.                         \tag{1}\label{eq:ex-cap-rank}
\]
Let $x_i\in X_i$, put $X_i^0=P_i x_i$, and set
$\lambda_i=(1+\|X_i^0\|)^{-1}$. Assume the following estimates hold
eventually in $I$:
\begin{align}
 \|L_i x_i\|&\le c,                                         \tag{2}\label{eq:ex-cap-tested}\\
 \|L_i(\lambda_i(x_i-X_i^0))\|&\le \frac{1}{2Z}.             \tag{3}\label{eq:ex-cap-remainder}
\end{align}
Then eventually
\[
                    \|X_i^0\|\le 2Zc+1.                    \tag{5}\label{eq:ex-cap-conclusion}
\]
\end{lemma}

\begin{proof}
Write $R_i=x_i-X_i^0$. By linearity and (2),
\[
 \|X_i^0\|\le Z\|L_iX_i^0\|
 \le Zc+Z\|L_iR_i\|.
\]
Since $\lambda_i=(1+\|X_i^0\|)^{-1}$, (3) implies
$Z\|L_iR_i\|\le(1+\|X_i^0\|)/2$. Therefore
\[
 \|X_i^0\|\le Zc+\tfrac12(1+\|X_i^0\|),
\]
and rearranging gives (5). All estimates were eventual, so their finite set
of thresholds can be replaced by their maximum.
\end{proof}

\begin{proposition}[One coefficient row is bounded at every period]
\label{prop:ex-fullmesh-cap-source}
Let $C\subseteq\prod_i X_i$ be one class of coefficient row families satisfying the following normalized-cost conditions. Suppose that the full family
$x=(x_i)$ belongs to $C$, that $C$ is closed under subtraction and
pointwise multiplication by scalar families whose absolute values are at most
one, and that every uniformly norm-bounded row family belongs to $C$. Suppose also that the detector estimates
\begin{enumerate}
\item $\|P_i z\|\le Z\|L_i(P_i z)\|$ for all $z\in X_i$ with a
      scale-independent $Z>0$;
\item for the primary paid trial there are assay values $b_i\in Y_i$ with
      $\|b_i\|\le v$ and $\|L_i x_i-b_i\|\to0$;
\item for every row family $z\in C$ satisfying $P_i z_i=0$ for every $i$,
      one has $\|L_i z_i\|\to0$.
\end{enumerate}
Then, eventually, uniformly over all indices in the directed mesh-period
system,
\[
                        \|P_i x_i\|\le 2Z(v+1)+1.
\]
The uniformity means that the estimates in (2) and (3) hold beyond one mesh
threshold for every period, as in the source full-mesh receipts.
\end{proposition}

\begin{proof}
From (2), eventually $\|L_i x_i\|\le v+1=:c$. Put
$X_i^0=P_i x_i$ and
\[
       \lambda_i=(1+\|X_i^0\|)^{-1},\qquad
       z_i=\lambda_i(x_i-X_i^0).
\]
The scalar family $\lambda=(\lambda_i)_i$ satisfies $|\lambda_i|\le1$, so
$\lambda x\in C$. The family $\lambda X^0$ is uniformly bounded by one, so
it belongs to $C$ as well. Therefore $z=\lambda(x-X^0)\in C$ by subtraction,
and $P_i z_i=0$ because
$P_i(x_i-P_ix_i)=0$. By (3),
for all sufficiently fine meshes $\|L_i z_i\|\le(2Z)^{-1}$. Since
$L_i z_i=\lambda_i L_i(x_i-X_i^0)$, estimate (1) gives
\[
 \|X_i^0\|\le Z\|L_iX_i^0\|
 \le Zc+Z\|L_i(x_i-X_i^0)\|
 \le Zc+\frac{1+\|X_i^0\|}{2}.
\]
Rearranging proves the claim. The common mesh threshold is the maximum of the
thresholds in (2) and (3).
\end{proof}

The cap alone bounds a projected row; the source separately uses that bound
to pay the difference trial. Here are the precise closure steps. Let $x$ be
the coefficient row for the requested insertion and let $F$ be its primary
paid trial. The exact equation gives an owed row
$e=\operatorname{upper}(x)-\operatorname{hops}(F)$, with a root-test cost
$u$ tending to zero. Applying the cap to $(x,F,e)$ bounds $Px$ uniformly.
Thus $y=Px$ is a uniformly bounded row family and its primary insertion is
paid by the bounded-coefficient receipt. The row $x-y$ lies in the normalized
cost class by the class closure, and $P(x-y)=0$. The sweep theorem therefore
makes $\operatorname{upper}(x-y)$ owed with a cost tending to zero. The
source identity is
\[
  \operatorname{hops}(F)
   =\operatorname{upper}(y)+\operatorname{upper}(x-y)-e.
\]
It follows from $\operatorname{upper}(x)=\operatorname{upper}(y)+
\operatorname{upper}(x-y)$ and the defining equation for $e$. The paid
budget is closed under finite addition and subtraction: once each vanishing
cost is below the unit allowance at a sufficiently fine mesh, the primary
paid budget together with the two unit owed budgets yields a finite paid
budget for the requested difference trial. This is exactly the projection-and-remainder argument;
it does not require the unsplit equation residual to be bounded by a quantity
that depends on $\|Px\|$.

\begin{theorem}[Full-mesh bounds for all finite differences]
\label{thm:ex-fullmesh-fixedbounds-producer}
Fix a family for which the preceding coefficient-row estimate and the
paid-remainder decomposition just described apply to every finite list of
fields and difference directions. Assume that bounded primary coefficient
rows have paid root-test budgets, that the exact periodic one-step difference
identity holds, and that translated insertions satisfy the product root-test
inequality in distinct interior cells.
Then for each finite field list $v$, difference list $l$, and $r>0$, there
are $B<\infty$ and a mesh threshold $N_0$ such that for every $N\ge N_0$,
every period $p$, and every periodic position tuple $x$ such that any two
distinct slot indices have sites separated by at least $r$ in physical units
(even if the corresponding insertions agree),
\[
                         |K_{N,p;v,l}(x)|\le B.              \tag{7}\label{eq:ex-fullmesh-bound}
\]
The threshold is independent of $p$; the bound is independent of both $N$
and $p$.
\end{theorem}

\begin{proof}
First fund one row. For a field label $w$, the constant coefficient row
$x_i=w$ has bounded normalized cost. The assumed primary receipt and the
preceding cap-and-remainder argument give a paid budget for each finite
difference list applied to that insertion, eventually at every mesh and
every period.
For a general finite list $v=(v_1,\ldots,v_n)$ and positional derivative
list, make one trial for each slot by applying to $v_j$ the subsequence of
directions assigned to that slot. The exact one-step identity is the
periodic discrete difference formula
\[
 F_j(x+ a_N e_\alpha)-F_j(x)
   =a_N\,\delta_{a_Ne_\alpha}F_j(x),
\]
with the normalization of the difference operator. Iterating this identity
and the exact product measure definition gives the desired multi-difference
kernel, including wrapped coordinates. The list is finite. Choose a common
centered scale beyond the finitely many thresholds in these paid records; at
that scale all its root-test budgets hold eventually in mesh, uniformly over
periods.

We now explain the two-scale tiling argument. Choose a tiling scale
$\ell$ and a finer recentering scale $k$ so that the physical tiling span
$a_\ell$ is below $r$ and
\[
                         6n a_k<a_\ell .
\]
First choose $\ell$ with $a_\ell<r$, then use $a_k\to0$ to choose $k$ with
the displayed margin. At a fine mesh $N$ beyond both scales, let $t$ be the
integer span of scale $k$ and let $T$ be a periodic tiling whose side is the
integer span of scale $\ell$. The aligned-scale identities give
$a_Nt=a_k$ and $a_N T.\mathrm{side}=a_\ell$. Thus
$n(2(3t))<T.\mathrm{side}$. The tiling placement lemma, applied with margin
$3t$, gives one translated tiling and an assignment of the $n$ insertion
sites to cells, each site lying at least $3t$ from the cell boundary.
This is uniform in the period: alignment gives the same physical side
$a_\ell$ for every period index, once $N$ is sufficiently fine.

The cell assignment is injective for a geometric reason. If two sites were
assigned to the same cell, their periodic maximum-coordinate separation
would be at most the integer tiling side $T.\mathrm{side}$ of that cell.
Multiplying by $a_N$ would give physical separation at most
$a_NT.\mathrm{side}=a_\ell<r$, contradicting the
assumption. The three-span margin is exactly what the centering lemma needs
to transport each local insertion from the centered tile to its assigned
cell.

For each occupied cell, recenter its insertion to the distinguished center
cell. The Wilson law is invariant under translations and the signed
coordinate symmetries used by the lattice. The centering lemma therefore
transports both the interior support condition and the same root-test budget
to the occupied cell. No statistical independence is asserted: all factors
are evaluated under the same Wilson law. Instead, the planted-cell
inequality applies to the translated copies in distinct cells and gives
\[
 \left|\mathbb E\prod_{j=1}^n F_j(x_j)\right|
       \leq \prod_{j=1}^n b_j,
\]
where $b_j$ is the paid root-test bound for the $j$th trial. Their finite
product is a constant independent of mesh and period. This proves
\eqref{eq:ex-fullmesh-bound}.
\end{proof}

\begin{remark}[What this estimate supplies]
The statement is the full-mesh separated-kernel estimate: each finite difference
order and each positive separation have one eventual mesh threshold valid at
all periods. The bounds at fixed period also control compact families of
separated positions and their coordinate ties. The estimate does not itself
give convergence of smeared correlations; that is the separate
same-history convergence estimate. With both inputs, the positional theorem
follows by the interpolation, compactness, and dominated-convergence argument
in Section~\ref{sec:ex-limits}.
\end{remark}
